\documentclass{article}
\usepackage[utf8]{inputenc}
\usepackage[spanish]{babel}
\usepackage{listings}
\usepackage{xcolor}
\usepackage{geometry}
\geometry{a4paper, margin=1in}

\% Configuración para código C++
\definecolor{codegreen}{rgb}{0,0.6,0}        
\definecolor{codegray}{rgb}{0.5,0.5,0.5}     
\definecolor{codepurple}{rgb}{0.58,0,0.82}   
\definecolor{backcolour}{rgb}{0.95,0.95,0.92}

\lstset{
    backgroundcolor=\color{backcolour},      
    commentstyle=\color{codegreen},
    keywordstyle=\color{magenta},
    numberstyle=\tiny\color{codegray},       
    stringstyle=\color{codepurple},
    basicstyle=\ttfamily\small,
    breakatwhitespace=false,
    breaklines=true,
    captionpos=b,
    keepspaces=true,
    numbers=left,
    numbersep=5pt,
    showspaces=false,
    showstringspaces=false,
    showtabs=false,
    tabsize=2,
    language=C++,
    inputencoding=utf8,
    extendedchars=true,
    literate={á}{{\'a}}1 {é}{{\'e}}1 {í}{{\'i}}1 {ó}{{\'o}}1 {ú}{{\'u}}1 {ñ}{{\~n}}1 {¿}{{\textquestiondown}}1 {¡}{{\textexclamdown}}1
}

\% Macros para las secciones de los apuntes
\newcommand{\quees}[1]{
    \section*{🤔 ¿QUÉ ES?}
    \hrulefill \\
    \#1
}

\newcommand{\dichodeotraforma}[1]{
    \section*{💬 DICHO DE OTRA FORMA}
    \hrulefill \\
    \#1
}

\newcommand{\diagrama}[1]{
    \section*{🎨 DIAGRAMA}
    \hrulefill \\
    \#1
}

\newcommand{\comofunciona}[1]{
    \section*{⚙️ ¿CÓMO FUNCIONA?}
    \hrulefill \\
    \#1
}

\newcommand{\complejidad}[1]{
    \section*{📱 COMPLEJIDAD (Big-O)}
    \hrulefill \\
    \#1
}

\newcommand{\entrevistas}[1]{
    \section*{🎯 EN ENTREVISTAS}
    \hrulefill \\
    \#1
}

\newcommand{\resumen}[1]{
    \section*{📋 RESUMEN RÁPIDO}
    \hrulefill \\
    \#1
}

\begin{document}

\title{BFS y DFS (Recorridos de Grafos)}
\author{Aldo Zepeda Montoya}
\date{\today}
\maketitle

\subsection*{CATEGORÍA: 03\_Grafos}

\quees{
Dos formas de explorar un mundo desconocido 🗺️:

1. BFS (Breadth-First Search): Como lanzar una piedra al agua 💧. 
   Ves las ondas expandirse nivel por nivel. Primero visitas a todos
   tus vecinos directos, luego a los vecinos de tus vecinos, y así 
   sucesivamente. ¡No dejas a nadie atrás en tu nivel!

2. DFS (Depth-First Search): Como explorar un laberinto 🌀. Sigues
   un solo camino hasta el fondo, hasta que ya no puedas avanzar. 
   Cuando llegas a un callejón sin salida, regresas (backtrack) y 
   pruebas el siguiente camino disponible.
}

\dichodeotraforma{
Son los dos algoritmos fundamentales para recorrer cualquier grafo 
o árbol. 

- BFS usa una Queue (FIFO) y garantiza encontrar el camino más corto
  en grafos no ponderados.
- DFS usa un Stack (LIFO) o Recursión. Es ideal para explorar todas 
  las posibilidades o detectar ciclos.

En ambos casos, es CRUCIAL usar una estructura `visited` (como un array
de bools o un set) para no entrar en bucles infinitos.
}

\diagrama{\begin{verbatim}
Grafo:
    [1]
   /   \
 [2]───[3]
  |     |
 [4]───[5]

BFS desde 1: 1 → 2,3 → 4,5
Orden: [1, 2, 3, 4, 5] (Nivel por nivel)

DFS desde 1: 1 → 2 → 4 → 5 → 3
Orden: [1, 2, 4, 5, 3] (Hasta el fondo)
\end{verbatim}}

\begin{lstlisting}[language=C++]
BFS (con Queue):
1. Metes el nodo inicial a la Queue y lo marcas como visitado.
2. Mientras la Queue no esté vacía:
   - Sacas el nodo del frente.
   - Metes a todos sus vecinos NO visitados a la Queue y los marcas.

DFS (con Recursión):
1. Marcas el nodo actual como visitado.
2. Para cada vecino NO visitado:
   - Llamas a la función DFS recursivamente con ese vecino.

📝 PSEUDOCÓDIGO
──────────────────────────────────────────────────────────────

// BFS iterativo
función BFS(inicio):
    q = crear Queue
    q.enqueue(inicio)
    marcar inicio como visitado
    MIENTRAS q NO vacía:
        u = q.dequeue()
        PARA cada vecino v de u:
            SI v NO visitado:
                q.enqueue(v)
                marcar v visitado

// DFS recursivo
función DFS(u):
    marcar u como visitado
    PARA cada vecino v de u:
        SI v NO visitado:
            DFS(v)

💻 CÓDIGO C++
──────────────────────────────────────────────────────────────

#include <iostream>
#include <vector>
#include <queue>
#include <stack>
using namespace std;

void bfs(int start, vector<vector<int>>& adj, vector<bool>& visited) {
    queue<int> q;
    q.push(start);
    visited[start] = true;

    while (!q.empty()) {
        int u = q.front(); q.pop();
        cout << u << " ";
        for (int v : adj[u]) {
            if (!visited[v]) {
                visited[v] = true;
                q.push(v);
            }
        }
    }
}

void dfs(int u, vector<vector<int>>& adj, vector<bool>& visited) {
    visited[u] = true;
    cout << u << " ";
    for (int v : adj[u]) {
        if (!visited[v]) dfs(v, adj, visited);
    }
}

int main() {
    int V = 6;
    vector<vector<int>> adj(V);
    // Definir grafo...
    
    vector<bool> visited(V, false);
    cout << "Recorrido BFS: ";
    bfs(0, adj, visited); // Suponiendo que empezamos en 0
    
    fill(visited.begin(), visited.end(), false);
    cout << "\nRecorrido DFS: ";
    dfs(0, adj, visited);
    
    return 0;
}

⏱️ COMPLEJIDAD (Big-O)
──────────────────────────────────────────────────────────────

  Algoritmo   │ Tiempo      │ Espacio
  ────────────┼─────────────┼──────────────
  BFS         │ O(V + E)    │ O(V)
  DFS         │ O(V + E)    │ O(V)

* V = Vértices, E = Edges.
* El espacio O(V) es para la Queue/Stack y el array `visited`.
\end{lstlisting}

\entrevistas{
✅ BFS para Shortest Path: Si te piden el camino más corto en un grafo
   SIN pesos (o una cuadrícula), BFS es la única opción correcta.

💡 DFS para Backtracking: Si el problema pide "encuentra todas las 
   rutas posibles" o "prueba todas las combinaciones", DFS es tu mejor
   amigo.

⚠️ Cycle Detection: 
   - Undirected: Si encuentras un vecino visitado que no es tu padre.
   - Directed: Necesitas 3 estados (Blanco, Gris, Negro). Si ves un 
     Gris, ¡hay ciclo!

🔑 Nivel de cada nodo: BFS es excelente para guardar la distancia 
   (nivel) de cada nodo respecto al origen.
}

\resumen{
• BFS = Queue = Nivel por nivel = Camino más corto.
• DFS = Stack/Recursión = Hasta el fondo = Exploración total.
• ¡Siempre usa un array `visited`!
• Ambos tardan O(V + E).
• Úsalos en matrices (grid), redes sociales, mapas y juegos.
════════════════════════════════════════════════════════════════
}

\end{document}